%%Property 3 alte Gleichung


\subsubsection*{Property 3: Fiscal Compensation Instruments}

Given that pandemic suppression is inherently contractionary (Property~2), the social planner requires compensating instruments to stabilize the economic state. Containment affects output through two distinct channels: it reduces the level of output by $\alpha_S$ per unit of lockdown intensity, and it erodes the economy's capacity to recover from the resulting gap by increasing output persistence through $\psi$.  A single fiscal instrument cannot address both channels simultaneously---closing the output gap through demand stimulus does not restore the recovery capacity that the lockdown has damaged, and preserving productive structures does not offset the contemporaneous output loss. This dual transmission motivates a disaggregation of fiscal policy into two instrument categories that differ in their design and their transmission to aggregate output.

Turning to Demand Injection (DI), denoted by $F_k^{DI}$, this category comprises measures that directly increase the flow of disposable income in the economy like stimulus checks (such as the U.S.\ Economic Impact Payments under the \textit{CARES Act}), enhanced unemployment benefits (such as the Federal Pandemic Unemployment Compensation of \$600 per week), direct transfers to households (such as the Australian \textit{JobSeeker} supplement or the Japanese universal cash transfer of
¥100,000 per resident), and tax reductions that raise disposable income (such as the German temporary VAT reduction from 19\% to 16\% in H2 2020). These instruments transmit through the standard multiplier channel, shifting the output gap by $\alpha_F^{DI} > 0$ per unit deployed. The magnitude of this multiplier is the central parameter governing the effectiveness of fiscal
compensation---in a pandemic no less than in normal times. Pre-pandemic estimates place expenditure multipliers in a range of approximately 0.5 to 1.5 \citep{ramey2018, blanchard2002}, while revenue-based multipliers range from near zero to cumulative effects of 2 to 3 depending on the identification strategy \citep{blanchard2002, romer2010}.

However, applying pre-pandemic multiplier estimates to a pandemic environment is questionable. When containment measures restrict the supply of goods and services, demand-side stimulus cannot translate into expenditure at the usual rate \citep{chetty2020, coibion2020, fariecastro2021, bayer2023c}. The effectiveness of $F_k^{DI}$ during containment therefore depends not only on the magnitude and design of the measure but appears to be strongly state-dependent. This state-dependency motivates a second category of instruments---Capacity Preservation (CP), denoted by $F_k^{CP}$--- to preserve the productive structure of the economy as without the resulting demand shortfall can exceed the initial supply disruption as shown by \citep{guerrieri2022}.

This category encompasses measures like short-time work schemes (such as the German \textit{Kurzarbeitergeld} or the UK Coronavirus Job Retention Scheme), payroll subsidies, government-guaranteed loans (such as the French \textit{Pr\^{e}ts Garantis par l'\'{E}tat} or the German KfW-Sonderprogramm). These measures help preserve the existing productive structure and thereby establish the foundation for a successful recovery and for effective demand-side stimulation. To see this effect and differs from DI, consider the role of $\rho_y$. In normal times, $\rho_y\in (0,1)$ captures the mean reversion of the output gap toward potential. During a pandemic, however, a prolonged lockdown threatens to alter this recovery dynamic. Firms with viable business cases but limited liquidity, that become insolvent during containment cannot resume production when restrictions lift. Employment relationships that dissolve must be rebuilt through costly search and matching \citep{Fujita2024}. Supply chains that break require time and resources to reconstitute. These mechanisms increase the persistence of the output gap---they effectively raise $\rho_y^{\text{}}$, slowing convergence back to potential \citep{Kozlowski2020, Cerra2020}. In the limit, $\rho_y^{\text{}} \to 1$ implies hysteresis: permanent output loss \citep{BLANCHARD1987}. Capacity preservation counteracts precisely this mechanism. \citet{deb2021} document a consistent pattern using high-frequency identification across 52 countries: liquidity-based measures (loans, guarantees) outperform income-support measures (transfers, tax relief) during periods of strict containment, while the ranking reverses upon reopening---with an average GDP multiplier of approximately 0.2 across both categories. Although their classification follows fiscal accounting conventions rather than the transmission-based distinction developed here, the state-dependent pattern aligns with the theoretical prediction: capacity preservation dominates when supply constraints bind, demand injection dominates when they do not. I formalize the effective persistence ($\rho_y^{\text{eff}}$) as:

%
\begin{equation}
    \rho_y^{eff}(S_k, F_k^{CP}) = \rho_y
      + \underbrace{\psi \, S_k}_{\substack{\text{lockdown-induced} \\ \text{hysteresis}}}
      - \underbrace{\eta \, F_k^{CP}}_{\substack{\text{capacity} \\ \text{preservation}}}
    \label{eq:rho_eff}
\end{equation}
%

where $\psi > 0$ governs the sensitivity of output persistence to lockdown intensity and $\eta > 0$ measures the effectiveness of capacity preservation in maintaining recovery speed. Both terms are multiplicative in the current output gap $y_k$: when the economy operates at potential ($y_k = 0$), neither hysteresis risk nor preservation benefit is operative---there is no output gap from which to recover, and no productive structures under stress.\footnote{The specification implies that $\rho_y^{eff}$ is determined by the \textit{contemporaneous} level of $S_k$ and $F_k^{CP}$. This means that structural damage ceases to accumulate once containment is lifted ($S_k = 0$), but does not capture the persistence of damage already incurred. For the pandemic application, where containment measures were adjusted gradually rather than terminated abruptly, this approximation is adequate. A generalization that tracks cumulative structural damage as a separate state variable would address this limitation at the cost of increased dimensionality.} When $\rho_y + \psi S_k \geq 1$ in the absence of accompanying capacity preservation, the output gap becomes explosive---the economy cannot recover without active structural support.This threshold provides a formal rationale for the universal combination of containment measures with capacity preservation programs observed during the pandemic. This existence is empirically covered by the fact that corporate insolvencies across OECD economies fell by 17\% on average in 2020 relative to 2019 despite the deepest recession since the Great Depression \citep{djankov2021, banerjee2021}. The classification of individual fiscal measures into the DI and CP categories, including instrument definitions, measurement conventions, and source documentation, is detailed in Section~\ref{sec:data}.
%Fusszeile auf eine Seite

Combining both channels, the complete output transition is formulated as:

%
\begin{equation}
    y_{k+1} = \underbrace{\left[\rho_y + \psi \, S_k - \eta \, F_k^{CP}\right]}_{\rho_y^{\text{eff}}} y_k
      \;-\; \alpha_S \, S_k
      \;+\; \alpha_F^{DI} \, F_k^{DI}
    \label{eq:output_full}
\end{equation}
%
The output gap evolves under three forces. The term $\rho_y^{\text{eff}} y_k$ propagates the inherited gap forward at an effective persistence rate from (\ref{eq:rho_eff}), the term $-\alpha_S S_k$ captures the direct contractionary impact of containment on current output, independent of the inherited gap\footnote{Since the pandemic originates as a contractionary shock and the suppression instrument reinforces this contraction, the output gap satisfies $y_k \leq 0$ along any relevant trajectory originating from the pre-pandemic steady state}. The term $+\alpha_F^{DI} F_k^{DI}$ represents the demand injection channel, which shifts the output level by a constant amount per unit deployed. The contemporaneous multiplier $\alpha_F^{DI}$ is state-independent; the cumulative impact of demand injection on future output, however, is state-dependent, as it propagates through $\rho_y^{\text{eff}}$ in subsequent periods---shaped by the containment and capacity preservation decisions that govern output persistence.

The partial derivatives of (\ref{eq:output_full}) with respect to each control reveal the state-dependent structure of the planner's problem and provide implications for the optimal composition of the policy mix (see Appendix~X for the full three-period forward iteration and the resulting optimality conditions and their intertemporality):

%
\begin{align}
    \frac{\partial y_{k+1}}{\partial S_k}       &= \psi \, y_k - \alpha_S < 0
    \label{eq:dy_dS} \\[6pt]
    \frac{\partial y_{k+1}}{\partial F_k^{DI}}  &= \alpha_F^{DI} > 0
    \label{eq:dy_dDI} \\[6pt]
    \frac{\partial y_{k+1}}{\partial F_k^{CP}}  &= -\eta \, y_k > 0
    \label{eq:dy_dCP}
\end{align}
%

\noindent First, the marginal output cost of containment, (\ref{eq:dy_dS}), is increasing in the depth of the recession. Each additional unit of lockdown is more costly when the economy is already weakened---not only because of the direct contractionary shock ($\alpha_S$), but because the lockdown further slows the recovery from an already deep output gap ($\psi |y_k|$). This convexity provides a formal rationale for early intervention: a lockdown imposed when $|y_k|$ is still small incurs lower marginal output costs than the same lockdown imposed after a prolonged contraction as already defined in the literature \citep{acemoglu2021}.

Second, the marginal output effect of demand injection ($\alpha_F^{DI}$) is constant, whereas the marginal effect of capacity preservation ($\eta |y_k|$) scales with the severity of the crisis. There exists a threshold $|y_k^*| = \alpha_F^{DI} / \eta$ beyond which the marginal output effect of capacity preservation exceeds that of demand injection. In a sufficiently deep recession, CP dominates DI on the output dimension \citep{deb2021}.

Third, capacity preservation reduces the marginal output cost of containment. When $\eta F_k^{CP} \approx \psi S_k$, the effective persistence returns to its baseline $\rho_y$ and the marginal cost of the lockdown reduces to $\alpha_S$---the pure contractionary shock without hysteresis amplification. CP does not raise output directly; it makes the lockdown economically sustainable.

These properties establish that demand injection and capacity preservation are complements, not substitutes. Moreover, the state-dependent structure implies a temporal sequencing logic: during strict containment ($S_k$ high, $|y_k|$ growing), the relative value of CP increases as it both counteracts hysteresis and reduces the output cost of the lockdown; as containment eases ($S_k$ declining), the relative value of DI rises as the demand channel operates more effectively. This complementarity provides a structural rationale for the universal adoption of mixed fiscal packages during the pandemic (see Section X).




















































\begin{itemize}
    \item \textbf{Responsiveness:} High $d_t$ triggers high $S_t$ (health priority), while high $y_t$ triggers high $F_t$ (economic priority).
    \item \textbf{Debt Constraint:} A high debt level $b_t$ dampens the usage of $F_t$ and may force a substitution towards $S_t$ if fiscal space is exhausted (demonstrating the Trilemma).
\end{itemize}


Further Steps:

- What does my findings say? Is the kind or the target or both relevant? because it needs a higher coordination
- Verify
- Solve t=0-> Decide on horizon
- Distinguish between Kind of F and/or Targets
- What does it mean? Interpretation
- Proof it empirically
\end{document}


%old version

\section*{1. Model Framework}
The model analyzes the decision problems of a social planner during a pandemic ($t=0,1,2$) who seeks to minimize three competing dimensions of damage. The planner has two instruments at their disposal: \textit{lockdown stringency} ($S_t \in [0,1]$) and \textit{fiscal support} ($F_t \ge 0$). The model ends in a terminal period ($t=3$) without controls.

\section*{2. The Social Planner's Objective Function}
The social planner minimizes expected discounted quadratic damages across all periods:
\begin{equation}
    \min_{\{S_t, F_t\}_{t=0}} \mathbb{E}_0 \left[ \sum_{t=0}^{2} \beta^t \left( w_y y_t^2 + w_d d_t^2 + w_b b_t^2 \right)\right] + \beta^3 \frac{1}{2} k_{b3} b_3^2
\end{equation}
where $y_t$ represents economic damage (output gap), $d_t$ health damage (excess mortality), and $b_t$ fiscal damage (debt accumulation). The parameters $w_i$ weight the respective policy objectives.

\section*{3. Dynamic Transition Equations}
The evolution of state variables follows a system of four linear difference equations under uncertainty:

\begin{align}
    \text{Economic Damage:} \quad & y_{t+1} = \rho_y y_t + \underbrace{\alpha_S S_t}_{\text{Lockdown Costs}} - \underbrace{\alpha_F F_t}_{\text{Fiscal Benefit}} \label{eq:econ} \\
    \text{Health Damage:} \quad & d_{t+1} = \rho_d d_t - \underbrace{\delta_S S_t}_{\text{Lockdown Benefit}} + \delta_\theta \theta_{t+1} \label{eq:health} \\
    \text{Fiscal Damage:} \quad & b_{t+1} = (1+r) b_t + \underbrace{\kappa_F F_t}_{\text{Fiscal Costs}} + \underbrace{\kappa_y y_t}_{\text{Auto. Stabilizers}} \label{eq:fiscal} \\
    \text{Pandemic Pressure:} \quad & \theta_{t+1} = \rho_\theta \theta_t - \underbrace{\phi_S S_t}_{\text{Lockdown Benefit}} + \epsilon_{t+1} \label{eq:pandemic}
\end{align}

\noindent \textbf{Legend:} $\theta_t$ describes the exogenous pandemic intensity. Stochastic shocks $\epsilon_{t+1}$ occur in the pandemic's evolution. Parameters $\rho_i \in (0,1)$ denote persistence; Greek letters ($\alpha, \delta, \kappa, \phi$) measure sensitivities.

\section*{4. The Mechanism of the Trilemma}
The model formalizes the fundamental trade-offs through the opposing signs of the impact parameters in equations \eqref{eq:econ}--\eqref{eq:fiscal}:

\begin{enumerate}
    \item \textbf{Health vs. Economy:} A stricter lockdown ($S_t \uparrow$) reduces health damage ($-\delta_S$ in Eq. \ref{eq:health}), but increases economic damage \textit{ceteris paribus} ($+\alpha_S$ in Eq. \ref{eq:econ}).
    \item \textbf{Economy vs. Fiscal Sustainability:} Higher fiscal support ($F_t \uparrow$) mitigates economic damage ($-\alpha_F$ in Eq. \ref{eq:econ}), but leads directly to an increase in public debt ($+\kappa_F$ in Eq. \ref{eq:fiscal}).
    \item \textbf{The Debt Feedback:} Insufficient use of $F_t$ leads to high economic damage $y_t$, which in turn increases the debt burden via automatic stabilizers (tax losses) ($+\kappa_y y_t$ in Eq. \ref{eq:fiscal}). This generates a non-trivial feedback loop.
\end{enumerate}

\noindent Since two instruments ($S, F$) act on three objectives ($y, d, b$), simultaneously minimizing all damages to zero is structurally impossible as long as $\theta_t > 0$.



\section*{5. Methodology}
The model is solved using \textit{backward induction}. Since the objective function is quadratic and the constraints are linear, the problem falls into the class of Linear-Quadratic (LQ) control problems. This ensures that the value functions are quadratic and the optimal policy rules are linear functions of the state vector $X_t = (y_t, d_t, b_t, \theta_t)$.

\section*{Period $t=3$ (Terminal Valuation)}
In the terminal period, no policy decisions are made ($S_3 = F_3 = 0$). The planner merely evaluates the residual states. The value function includes the standard weights plus the additional terminal debt penalty parameter $k_{b3}$:
\begin{equation}
    V_3(y_3, d_3, b_3) = \beta^3 \left[ w_y y_3^2 + w_d d_3^2 + \tilde{w}_b^{(3)} b_3^2 \right]
\end{equation}
where the effective terminal debt weight is defined as $\tilde{w}_b^{(3)} = w_b + \frac{1}{2}k_{b3}$. This step establishes the terminal boundary condition for the optimization.

\section*{Period $t=2$ (Deterministic Optimization)}
In $t=2$, the pandemic is assumed to subside ($\epsilon_3 = 0$), making the transition to $t=3$ deterministic. The planner solves:
\begin{equation}
    V_2(X_2) = \min_{S_2, F_2} \left\{ \beta^2 \left( w_y y_2^2 + w_d d_2^2 + w_b b_2^2 \right) + V_3(X_3) \right\}
\end{equation}
Subject to the transition equations for $X_3$. The FOCs yield a system of linear equations. The optimal policies are functions of the current state:
\begin{align}
    S_2^*(X_2) &= \phi_{S2}^y y_2 + \phi_{S2}^d d_2 + \phi_{S2}^b b_2 + \phi_{S2}^\theta \theta_2 \\
    F_2^*(X_2) &= \phi_{F2}^y y_2 + \phi_{F2}^d d_2 + \phi_{F2}^b b_2 + \phi_{F2}^\theta \theta_2
\end{align}
These policy functions quantify how the planner reacts to late-stage damages (e.g., austerity measures if $b_2$ is high).

\section*{Period $t=1$ (Stochastic Optimization)}
In $t=1$, the planner faces uncertainty regarding the pandemic evolution in the next period ($\theta_2 = \rho_\theta \theta_1 - \phi_S S_1 + \epsilon_2$). The Bellman equation incorporates expectations:
\begin{equation}
    V_1(X_1) = \min_{S_1, F_1} \left\{ \beta \left( w_y y_1^2 + w_d d_1^2 + w_b b_1^2 \right) + \beta \mathbb{E}_1 \left[ V_2(X_2) \right] \right\}
\end{equation}
Due to the \textit{Certainty Equivalence} property of LQ models, the optimal policy depends only on the expected future state $\mathbb{E}[\theta_2]$ and not on the variance $\sigma_2^2$ (see proof). The solution procedure yields the policy rules $S_1^*(X_1)$ and $F_1^*(X_1)$, which now internalize the dynamic continuation value $V_2$. The derived policy coefficients are a complex system (\ref{eq:matrix}) reveal the higher need for a coordinated structural response of the planner under uncertainty.

\begin{equation}
    \begin{pmatrix} 
    S_1^* \\ 
    F_1^* \end{pmatrix} 
    = Q^{-1} N_1 
    \begin{pmatrix} 
    y_1 \\ 
    d_1 \\ 
    b_1 \\ 
    \theta_1 
    \end{pmatrix} 
    = A_1 (\text{Policy Matrix}) 
    = 
    \begin{pmatrix} 
    a_{1S}^y & a_{1S}^d & a_{1S}^b & a_{1S}^\theta \\ 
    a_{1F}^y & a_{1F}^d & a_{1F}^b & a_{1F}^\theta 
    \end{pmatrix}
    \label{eq:matrix}
\end{equation}

%OLD


\section{Problem Statement}
We consider a stochastic linear system:
\begin{equation}
x_{k+1} = A_k x_k + B_k u_k + w_k, \quad k=0, 1, \dots, N-1
\end{equation}
where $E[w_k] = 0$ and $w_k$ is independent of $x_k$ and $u_k$. The objective is to minimize:
\begin{equation}
J = E \left\{ x_N' Q_N x_N + \sum_{k=0}^{N-1} (x_k' Q_k x_k + u_k' R_k u_k) \right\}
\end{equation}

\section{Dynamic Programming Algorithm}

\subsection{Step 1: Initialization at $k=N$}
The optimal cost at the final stage depends only on the terminal state:
\begin{equation}
J_N(x_N) = x_N' Q_N x_N \implies P_N = Q_N
\end{equation}

\subsection{Step 2: The Recursive Step}
For $k = N-1, N-2, \dots, 0$, the Bellman equation is:
\begin{equation}
J_k(x_k) = \min_{u_k} E_{w_k} \left\{ x_k' Q_k x_k + u_k' R_k u_k + J_{k+1}(Ax_k + Bu_k + w_k) \right\}
\end{equation}
Assuming $J_{k+1}(x) = x' P_{k+1} x + \text{constant}$, and substituting the dynamics:
\begin{equation}
J_k(x_k) = x_k' Q_k x_k + \min_{u_k} \left\{ u_k' R_k u_k + (Ax_k + Bu_k)' P_{k+1} (Ax_k + Bu_k) + E[w_k' P_{k+1} w_k] \right\}
\end{equation}

\subsection{Step 3: Minimization and Policy Gain}
Differentiating the expression inside the brackets with respect to $u_k$ and setting it to zero:
\begin{equation}
2 R_k u_k + 2 B' P_{k+1} (Ax_k + Bu_k) = 0
\end{equation}
Solving for $u_k$ yields the optimal control law:
\begin{equation}
u_k^* = -K_k x_k, \quad \text{where } K_k = (B' P_{k+1} B + R_k)^{-1} B' P_{k+1} A
\end{equation}

\subsection{Step 4: The Riccati Equation}
Substituting $u_k^*$ back into the cost function, we obtain the recursive update for the matrix $P_k$:
\begin{equation}
P_k = A' P_{k+1} A - (A' P_{k+1} B) (B' P_{k+1} B + R_k)^{-1} (B' P_{k+1} A) + Q_k
\end{equation}

\section{The Role of the Disturbance $w_k$}
In this model, the noise term $w_k = (0, \delta_w \varepsilon_k, 0, \varepsilon_k)'$ affects the state transitions. However, because the system is linear and the cost is quadratic, the \textbf{Certainty Equivalence Principle} holds:
\begin{enumerate}
    \item The optimal feedback gain $K_k$ is \textit{independent} of the noise variance.
    \item The gain $K_k$ is the same as in the deterministic case ($w_k=0$).
    \item The expected cost increases by $\sum_{k=0}^{N-1} E[w_k' P_{k+1} w_k]$ due to the stochastic nature of the pandemic.
\end{enumerate}



%5Alte Reihenfolge-> Research Question danach


\subsection{Research Question}

%Quellen hinein!!!Besser schreiben, Joint Optimization, Ergänzung zu den anderen Modellen und auch Parameter in this Sample, Data, für die Zukunft, vlt auch auf Natur

%Fiscal, Distinction, Endogenous fiscal feedback, Joint optimization, Cross Country Empirical Estimation delivers optimal Paths and Vergleiche


The constitutive properties established above define the pandemic environment as a constrained optimization problem with multiple unstable state dimensions, structurally coupled instruments, and circular fiscal feedback. A growing literature has formalized subsets of this problem. Alvarez, Argente, and Lippi (2021) solve for the optimal lockdown trajectory of a planner trading off fatalities against output costs, but treat the economy as a linear production function with no fiscal dimension. Eichenbaum, Rebelo, and Trabandt (2021) endogenize the behavioral recession by embedding SIR dynamics
into a macroeconomic equilibrium, yet the government's only instrument is a Pigouvian containment tax---no fiscal compensation or debt dynamics enter the problem. Kaplan, Moll, and Violante (2020) come closest to the full trilemma: they integrate SIR dynamics into a heterogeneous-agent model with both lockdown and fiscal transfers, constructing a \textit{pandemic possibility frontier} that maps policy combinations to
the joint distribution of mortality and economic welfare losses. Their analysis establishes the binding health--economy trade-off and demonstrates that fiscal transfers reshape its distributional incidence. However, government debt enters as a residual accounting identity, not as an optimization target, and fiscal policy is modeled as undifferentiated lump-sum transfers rather than as a composition of instruments with distinct transmission mechanisms.\footnote{For a systematic comparison of modeling
approaches in the pandemic-macro literature, see Table~?? in Appendix~??.}

Three gaps emerge from this body of work. First, no existing framework incorporates fiscal sustainability as a co-equal optimization target alongside health protection and output stabilization---the problem remains a dilemma rather than a trilemma. Second, the fiscal instruments available to the planner have not been disaggregated by their transmission mechanism: demand injection, which operates as a level shift on the output gap ($\alpha_F^{DI}$), and capacity preservation, which reduces the persistence of
the contraction ($\eta$), generate qualitatively different dynamic profiles and carry different fiscal costs ($\kappa_F^{DI}$ vs.\ $\kappa_F^{CP}$), yet existing models either omit fiscal policy entirely or treat it as a single aggregate instrument. Third, the instrument-specific structural parameters that govern the composition problem have not been estimated jointly under pandemic conditions using a
transmission-mechanism-based disaggregation---existing cross-country evidence (Deb et al., 2021) benchmarks aggregate fiscal effectiveness but not compositional optimality.

Closing these gaps requires both a theoretical framework that jointly optimizes over containment and disaggregated fiscal instruments under an explicit debt constraint, and an empirical basis that provides the instrument-level granularity to identify the relevant structural parameters. This paper contributes both.

The research question is therefore:

\medskip
\noindent\textit{Given a country's pandemic trajectory, which fiscal composition---across demand injection and capacity preservation, at what intensity, and at which point in the containment trajectory---jointly minimizes output loss and debt accumulation? And what are the costs---in foregone stabilization and excess
debt---of deviating from this composition?}



To analyze this question, this paper makes three contributions.

First, I develop a formal optimization framework---an iterative Linear-Quadratic
Regulator (iLQR) following \citet{bertsekas2017}---that solves the full trilemma: the
planner jointly optimizes containment intensity $S_k$ and the fiscal instruments
$F_k^{DI}$ and $F_k^{CP}$ across the multi-dimensional state space
$(y_k, d_k, b_k, \theta_k)$, capturing the bilinear interactions between containment
and economic persistence that define the pandemic trilemma. Existing approaches solve
simpler variants of this problem: Alvarez et al.\ (2021) apply Pontryagin's maximum
principle to a single-instrument continuous-time problem, and Eichenbaum et al.\
(2021) solve a Ramsey problem over a scalar containment tax. The iLQR extends this
to a multi-instrument setting while preserving analytical tractability: the
backward Riccati recursion yields time-varying feedback gains $K_k$ that are directly
interpretable as the planner's marginal policy response to state deviations, without
requiring numerical value function iteration over a high-dimensional state
space.\footnote{The bilinear term $\phi_S S_k \theta_k$ renders the full problem
non-convex. The iterative linearization converges to a local optimum. In the
economically relevant region---where the planner responds monotonically to rising
infections---local and global solutions coincide, as verified numerically in
Appendix~\ref{app:convergence}.} The resulting state-dependent feedback law
$u_k^* = -K_k x_k$ yields closed-form optimality conditions
(Appendix~\ref{app:optimality}) that characterize the optimal policy
trajectory---including the optimal intensity, timing, and composition of the fiscal
response---as a function of the evolving state of the system. The quadratic objective
penalizing $(y_k, d_k, b_k)$ can be interpreted as a second-order approximation to a
welfare function in which the planner assigns explicit weights to output stabilization,
mortality reduction, and fiscal sustainability---a Tinbergen-style formulation that
takes the multi-objective nature of pandemic policy as given rather than deriving it
from a single representative agent's preferences.\footnote{This abstracts from the
distributional dimension emphasized by Kaplan, Moll, and Violante (2020). The
present framework addresses a logically prior question: what is the optimal
\textit{aggregate} fiscal composition? The distributional incidence of a given
composition is a complementary but distinct problem.}

Second, I construct a new micro-level database of 2,290 fiscal policy measures across
all 38 OECD member states during 2020--2021 \citep{pesenti2025}. Unlike existing
databases that aggregate fiscal responses at the package level, this dataset
disaggregates individual measures by instrument type, target group, and IMF fiscal
accounting category---the granularity required to operationalize the distinction
between demand injection and capacity preservation that is central to the theoretical
framework. Each of the 49 policy instrument codes is mapped to the theoretical
categories (DI or CP) based on a transparent classification protocol that identifies
the primary transmission channel: instruments that sustain aggregate expenditure
through direct income replacement are classified as demand injection, while
instruments that preserve productive capacity by maintaining employment relationships,
firm solvency, or sectoral viability are classified as capacity
preservation.\footnote{The classification is documented in detail in Section~??. A
subset of instruments---notably wage subsidies---operates through both channels
simultaneously. These are allocated based on the dominant transmission mechanism, and
all quantitative results are subjected to robustness checks under alternative
classification rules (Section~??).} Without this disaggregation, the composition
question cannot be empirically addressed: aggregate fiscal impulse measures conflate
instruments with fundamentally different transmission mechanisms and fiscal cost
profiles.

Third, I estimate the structural parameters of the transition system---in particular
$\alpha_F^{DI}$, $\eta$, $\kappa_F^{DI}$, $\kappa_F^{CP}$, $\phi_S$, and
$c_H$---using panel Local Projections \citep{jorda2005, jorda2025} across 38 OECD
economies during 2020--2021. The identification strategy exploits cross-country
variation in fiscal composition conditional on comparable containment trajectories. A
two-stage approach first estimates the policy reaction function---fiscal composition
as a function of containment stringency, pre-crisis fiscal position, and institutional
characteristics---and then uses deviations from predicted composition as identifying
variation.\footnote{The exclusion restriction requires that residual variation in
fiscal composition affects economic outcomes only through the fiscal channel.
Section~?? discusses potential threats, including the possibility that political
economy variables driving composition choices (government ideology, coalition
structure) independently affect expectations and hence outcomes.} The estimated
parameters calibrate the theoretical model and enable two counterfactual exercises.
The primary exercise takes each country's observed containment path as given and
computes the conditionally optimal fiscal composition---quantifying the excess debt
attributable to suboptimal instrument choice, holding the containment trajectory
constant. This conditional design reflects the institutional reality that fiscal
authorities responded to containment decisions rather than determining them, and
isolates the margin over which fiscal policymakers had actual discretion. A secondary
exercise computes the fully optimal trajectory across all three instruments, providing
an upper bound on the welfare gains from joint optimization of containment and fiscal
policy.

This question is not only of historical interest. The fiscal legacy of the
pandemic---average OECD government debt increased by approximately XX percentage
points of GDP between 2019 and 2021---directly shaped the fiscal space available when
subsequent shocks materialized. The inflation surge of 2022--2023 raised sovereign
borrowing costs across advanced economies \citep{jorda2023}; the energy crisis
following the invasion of Ukraine required additional fiscal responses in an
environment of already elevated debt; and the prevailing macroeconomic uncertainty
continues to constrain fiscal policy options. Countries that achieved comparable health
and economic outcomes during the pandemic but at lower fiscal cost---through more
effective instrument composition---entered this sequence of subsequent shocks in a
structurally stronger position. Understanding which fiscal strategies achieved this
outcome, and why, is therefore directly relevant to the current fiscal policy
environment and to the management of future crises that share the pandemic's
structural characteristics.




























\subsubsection{Contribution}
To analyze this research question, this paper makes three contributions.

First, I develop a formal optimization framework---an iterative Linear-Quadratic Regulator (iLQR) following \citet{bertsekas2017}---that solves the full trilemma: the planner jointly optimizes containment
intensity $S_k$ and the fiscal instruments $F_k^{DI}$ and $F_k^{CP}$ across the multi-dimensional state space ($y_k$, $d_k$, $b_k$, $\theta_k$), capturing the bilinear interactions between containment and economic persistence that define the pandemic trilemma. The resulting state-dependent feedback law $u_k^* = -K_k x_k$ yields closed-form optimality conditions (Appendix~\ref{app:optimality}) that characterize the optimal policy trajectory---including the optimal intensity, timing,
and composition of the fiscal response---as a function of the evolving state of the system.

Second, I construct a new micro-level database of 2,290 fiscal policy measures across all 38 OECD member states during 2020--2021 \citep{pesenti2025}. Unlike existing databases that aggregate fiscal responses at the package level, this dataset disaggregates individual measures by instrument type, target group, and IMF fiscal accounting category---the granularity required to operationalize the distinction between demand injection and capacity preservation that is central to the
theoretical framework. Without this disaggregation, the composition question cannot be empirically addressed: aggregate fiscal impulse measures conflate instruments with fundamentally different transmission mechanisms and fiscal cost profiles.

Third, I estimate the structural parameters of the transition system---in particular $\alpha_F^{DI}$, $\eta$, $\kappa_F^{DI}$, $\kappa_F^{CP}$, $\phi_S$, and $c_H$---using panel Local Projections \citep{jorda2005, jorda2025} across 38 OECD economies during 2020--2021. The identification strategy exploits policy reaction function residuals to isolate exogenous variation in both containment stringency and fiscal support. The estimated parameters calibrate the theoretical model and enable two counterfactual exercises. The primary exercise takes each country's observed containment path as given and computes the conditionally optimal fiscal composition---quantifying the excess debt attributable to suboptimal instrument choice, holding the containment trajectory constant. This conditional design reflects the institutional
reality that fiscal authorities responded to containment decisions rather than determining them, and isolates the margin over which fiscal policymakers had actual discretion. A secondary exercise computes the fully optimal trajectory across all three instruments, providing an upper bound on the welfare gains from joint optimization of containment and fiscal policy.

This question is not only of historical interest. The fiscal legacy of the pandemic---average OECD government debt increased by approximately XX percentage points of GDP between 2019 and 2021---directly shaped the fiscal space available when subsequent shocks materialized. The inflation surge of 2022--2023 raised sovereign borrowing costs across advanced economies \citep{jorda2023}; the energy crisis following the invasion of Ukraine required additional fiscal responses in an environment of already elevated debt; and the prevailing macroeconomic uncertainty continues to constrain fiscal policy options. Countries that achieved comparable health and economic outcomes during the pandemic but at lower fiscal
cost---through more effective instrument composition---entered this sequence of subsequent shocks in a structurally stronger position. Understanding which fiscal strategies achieved this outcome, and why, is therefore directly relevant to the current fiscal policy environment and to the management of future crises that share the pandemic's structural characteristics.

Cutler and Summers (2020) estimate the total economic cost of the pandemic in the
United States alone at \$16 trillion, approximately 90\% of annual GDP, underscoring the
scale of the fiscal challenge \citep{CutlerSummers2020}. Das sagt nur die Erwartungen -> Facts einfügen




\subsubsection*{Roadmap}

The remainder of the paper proceeds as follows. Section~\ref{sec:model}
formalizes the complete transition system and the planner's optimization problem
within the iLQR framework, derives the optimality conditions, and characterizes
the state-dependent feedback law. Section~\ref{sec:data} describes the fiscal
policy database and the construction of the instrument-category variables.
Section~\ref{sec:empirics} presents the empirical strategy, including the policy
reaction functions and the Local Projection specification.
Section~\ref{sec:results} reports the estimation results, maps the estimated
coefficients to the structural parameters of the model, computes optimal policy
trajectories, and quantifies the deviation costs of observed fiscal strategies
relative to the model benchmark across OECD economies.
Section~\ref{sec:conclusion} discusses implications for future crisis management
and concludes.

I estimate impulse response functions using panel Local Projections
\citep{jorda2005, jorda2025} at horizons $h = 0, 1, \ldots, H$ for each
outcome variable $x \in \{y, d, b, \theta\}$. The contemporaneous
coefficients ($h = 0$) identify the structural parameters of the transition
system---$\alpha_S$, $\alpha_F^{DI}$, $\phi_S$, $\kappa_F^{DI}$,
$\kappa_F^{CP}$, and $c_H$---while the dynamic profile across horizons
identifies the persistence-related parameters $\eta$ and $\psi$ through
the horizon-dependent attenuation of containment and fiscal shocks. The
full set of estimated impulse responses additionally serves as a
validation benchmark: if the structural parameters are correctly
identified, forward simulation of the calibrated model should approximate
the LP estimates across horizons.
%Notes
The LP specification identifies the dynamic effect of a fiscal impulse at
time $t$ on outcomes at horizon $t+h$. Since governments implemented
sequential fiscal packages throughout the pandemic, the estimated impulse
response reflects the effect of a fiscal shock in an environment where
subsequent policy adjustments also occur---the ``typical policy sequence''
interpretation emphasized by \citet{jorda2025}. This interpretation is
consistent with the theoretical model, where the planner selects a new
policy vector $u_k$ in every period. As a robustness check, I include
leads of the fiscal shock variable to control explicitly for subsequent
interventions (Appendix~\ref{app:robustness}); the baseline results are
not sensitive to this augmentation.


\newpage

%%%ALT
\paragraph{Summary}


Auch kurz verweisen auf die Vergleiche zu anderen Krisen (Appendix), die Einzigartigkeit zeigen....
(Some aktuelle Quellen, und auch noch die Energierise, di eUnsicherheit heutzutage).

So here wrap up warum diese theoretisch implikation die Grundlage legt für eine wichtige Forschung die zu der Literatur beiträgt, hilft für diezukunft und auch einen beitrag zu erklärung der heutigen situation liefert und wie dies verhindert werden hätte können. Wer war am besten?? Nur mein Datensatz erlaubt es diese Unterscheidung sauber zu machen und ermöglicht diese Analyse überhaupt.-> Maybe all das in das Summary und dort direkt die Research Question staten, dann die Lösung des Modells inkl. Literatur review (parameter kalibrieren und maybe comparative statics, die Daten und Schätzung, Ergebnisse und Counterfactuals.

More Research Motivation... modell zeigt die theoretisch ewirkung, es ist also intertemporal, state dependend und mix macht es aus, deshalb das herausfinden weil....This question is important to study the optimal behavior in future crises but also brings up som eimportant implications about the current debt situation (QUELLE). 
Diese 4 Properties zeigen somit, das während COVID ein Trilemma bestand (Kausalkette kurz einfügen).  Diese Situation unterschied sich stark von anderen Krisen und ist in diesem Umfang aufgrund des persistenten Gesundheitsshock einzigartig. Vergleicht man die Situation mit anderen Krisen wird sichtbar das nur di eProblematik mit der Kimakrise am nähesten an diese Situation kommt. In einer Klimakrise gibt es auch einen persistenten Schock (Property 1), ein Instrument das Einzudämmen (Property 2), ein Instrument die negativen Folgen abzuschwächen (Property 3) und eine autmatische Dynamik ohne Eingriff. Jedoch ist der Zeithorizont ein deutlich andere und exogene wie auch endogene Faktoren wie technologischer Fortschritt können dieses Trilemma auflösen (siehe Vergleich Imfpung).

Furthermore, looking at the debt transition, we can observe that the unconditionally explosive eigenvalue $(1+r) > 1$ implies that debt accumulated during the pandemic compounds beyond the planning horizon regardless of post-crisis policy. During the crisis itself, the feedback from debt to output is negligible over
the short pandemic horizon and is therefore excluded from this specification.\footnote{The standard channels through which debt affects output---sovereign
risk premia, crowding out of private investment, anticipation of future fiscal consolidation---operate over multi-year horizons and are unlikely to materially alter
the planner's quarterly decisions during an acute crisis.} 
However, the planner's debt minimization incentive does not derive from within-crisis feedback but from the exposure that terminal debt $b_N$ creates to post-pandemic conditions.

The post-pandemic environment materialized precisely the risk that motivates this concern. The inflation surge of 2022--2023 prompted sharp increases in policy rates
across advanced economies \citep{jorda2023}, raising the effective $r$ on newly issued and refinanced sovereign debt. Countries that exited the pandemic with higher $b_N$
faced larger absolute increases in debt servicing costs, tighter fiscal space for subsequent shocks, and---in some cases---elevated sovereign spreads reflecting market
reassessment of fiscal sustainability. This observation reinforces the planner's incentive to minimize the debt footprint of pandemic interventions: not because debt
constrains policy during the crisis, but because it determines the economy's vulnerability to the post-pandemic interest rate environment. The terminal cost $W_b$ in the objective function (Section~X) formalizes this concern by imposing a penalty on $b_N$ that disciplines the planner's intertemporal trade-off between current
stabilization and future fiscal space. 





The preceding analysis establishes that the effectiveness of fiscal compensation
depends jointly on three dimensions: the \textit{composition} of the fiscal package
($F_k^{DI}$ versus $F_k^{CP}$), the \textit{state} of the economy at the time of
deployment ($y_k$, $\theta_k$), and the \textit{timing} relative to the containment
trajectory ($S_k$). The output benefit of demand injection is constant
($\alpha_F^{DI}$), whereas the benefit of capacity preservation scales with the depth
of the recession ($\eta |y_k|$). The fiscal cost of demand injection is largely
predetermined ($\kappa_F^{DI} \approx 1$), whereas the cost of capacity preservation
is contingent on outcomes that the planner simultaneously influences. And both
instruments interact with containment: a stricter lockdown increases the required
fiscal volume while compressing the multiplier through which fiscal support transmits
to output.

The existing literature has not addressed these interdependencies in a unified
framework. Studies of pandemic fiscal multipliers treat the composition of fiscal
packages as given \citep{auerbach2022,bayer2023}; analyses of containment policy
abstract from the fiscal dimension \citep{acemoglu2021,alvarez2021}; and the
interaction between containment stringency and fiscal effectiveness---though documented
in isolated settings \citep{XXauerbach_defense_spendingXX}---has not been integrated
into an optimization framework that accounts for the intertemporal debt consequences
of instrument choice. The question of which fiscal composition minimizes the debt
footprint of a given level of output stabilization, conditional on the prevailing
containment intensity and the economy's structural characteristics, therefore remains
open. This is the question that the model framework developed in Section~X is
designed to answer.

\paragraph{Post-pandemic debt sustainability.}

Beyond the optimal management of the crisis itself, the composition question carries
direct implications for the fiscal position that countries inherit after the pandemic.
The debt transition (\ref{eq:debt_transition}) contains an unconditionally explosive
eigenvalue: $(1+r) > 1$ implies that any debt accumulated during the crisis compounds
beyond the planning horizon regardless of post-crisis policy. During the pandemic
itself, the feedback from debt to output is negligible over the short crisis horizon
and is therefore excluded from this specification.\footnote{The standard channels
through which public debt affects output---sovereign risk premia, crowding out of
private investment, anticipation of future fiscal consolidation---operate over
multi-year horizons and are unlikely to materially alter the planner's quarterly
decisions during an acute crisis lasting six to eight quarters.}

The planner's incentive to minimize debt does not derive from this within-crisis
feedback but from the \textit{exposure} that terminal debt $b_N$ creates to
post-pandemic conditions that are unknown at the time of the fiscal deployment
decision. The post-pandemic environment materialized precisely the risk that motivates
this concern: the inflation surge of 2022--2023 prompted sharp increases in policy
rates across advanced economies \citep{jorda2023}, raising the effective $r$ on newly
issued and refinanced sovereign debt. Countries that exited the pandemic with higher
$b_N$ faced three compounding disadvantages: larger absolute increases in debt
servicing costs, reduced fiscal space for subsequent shocks, and---in some
cases---elevated sovereign spreads reflecting market reassessment of fiscal
sustainability. A fiscal strategy that achieved equivalent output stabilization during
the crisis at lower terminal debt $b_N$ would have yielded a measurable advantage in
this post-pandemic environment---an advantage that is invisible in analyses confined to
the crisis period. The terminal cost $W_b$ in the objective function (Section~X)
formalizes this concern: it imposes a penalty on $b_N$ that disciplines the planner's
intertemporal trade-off between current stabilization and future fiscal vulnerability.


%Entwurf
\subsection*{Summary}

The four constitutive properties jointly establish a circular interdependence that precludes the simultaneous attainment of the planner's three objectives. Property~1 renders suppression necessary: unchecked viral dynamics ($\rho_\theta > 1$) generate escalating mortality. Property~2 binds suppression to economic contraction: any intervention that reduces transmission simultaneously influences output, as both effects operate through the same channel of contact restriction ($\alpha_S, \phi_S > 0$). Property~3 introduces the fiscal compensation instrument and its inescapable debt consequence: stabilizing the economy requires fiscal support ($\alpha_F^j > 0$), but every instrument category carries a strictly positive fiscal cost ($\kappa_F^j > 0$). Property~4 closes the circle by showing that even fiscal abstinence fails to preserve the budget: through automatic stabilizers ($\gamma_y > 0$), the output contraction induced by suppression erodes the fiscal position independently of any discretionary spending decision. The planner is therefore trapped: protecting health destabilizes the economy, compensating the economy generates debt, and inaction accumulates debt passively---no feasible policy combination restores all three state variables to the origin simultaneously.

This structure is inherently dynamic and cannot be resolved by static optimization. The suppression instrument operates with a two-period lag ($S_k \to \theta_{k+1} \to d_{k+2}$), whereas fiscal compensation transmits within one period ($F_k \to y_{k+1}$). Combined with the explosive epidemiological dynamic ($\rho_\theta > 1$), this temporal asymmetry implies that delayed intervention is convexly costly: each period of inaction widens the gap that future policy must close under increasingly unfavorable conditions. The planner must therefore optimize intertemporally, weighing current costs against future state trajectories across all three dimensions.

This configuration is structurally unique to pandemics. While other crisis types---wars, natural disasters, financial crises, climate change---share individual properties, none combines all four simultaneously. Wars destroy physical capital but lack the endogenous, self-reinforcing growth dynamic of $\rho_\theta > 1$; financial crises are endogenous and do not require a suppression instrument with inherently contractionary side effects; natural disasters are impulse shocks ($\rho < 1$) that dissipate without sustained intervention. Section~X applies the four properties as a diagnostic framework to establish this uniqueness formally.

Two central questions emerge from the trilemma structure that the existing literature has not addressed in a unified framework. First, the instrument-specific fiscal multipliers $\alpha_F^j$ and their associated debt costs $\kappa_F^j$ have not been estimated jointly across the three categories---expenditure, revenue, and credit support---under pandemic conditions, where supply constraints and containment stringency fundamentally alter the transmission of fiscal impulses. Second, the optimal composition of the fiscal package $F_k = (F_k^{Exp}, F_k^{Rev}, F_k^{L\&G})'$ conditional on the prevailing level of suppression $S_k$ and the economy's structural characteristics remains an open problem. Addressing these questions requires a framework that captures the simultaneous management of multiple interconnected state dimensions, the feedback loops between health, economic, and fiscal outcomes, and the intertemporal consequences of policy choices under an explosive epidemiological dynamic.

This motivates the iterative Linear-Quadratic Regulator (iLQR) framework developed in Section~X. The LQR approach is suited to this problem for three reasons. First, it accommodates the multi-dimensional state space $(y_k, d_k, b_k, \theta_k)$ and the multi-instrument control vector $(S_k, F_k^{Exp}, F_k^{Rev}, F_k^{L\&G})$ within a unified optimization, capturing the cross-dimensional trade-offs that define the trilemma. Second, the quadratic cost structure reflects the planner's disproportionate aversion to large deviations---healthcare system collapse, sovereign default, economic depression---and penalizes distance from the pre-pandemic equilibrium at the origin. Third, the iterative extension (iLQR) accommodates the bilinear suppression dynamics established in Property~2, where the multiplicative interaction $\rho_\theta(1 - \phi_S S_k)\theta_k$ breaks the linearity required by standard LQR but can be handled through successive linearization around the current trajectory. The resulting state-dependent feedback law $u_k^* = -K_k x_k$ provides interpretable policy rules that map directly onto the structural parameters identified by the four properties.


%Fazit aufzeigen das dies auch noch dynamisch ist und alles zusammenhängt: Was bedeutet das also für den Planner und was sind offene Fragen die noch bestehen in der Forschung. Auch uniqness hier betonen und auf die mathematische formulierung überleiten.

\newpage
\vspace{5mm}

%% --- Dynamics ---


%%Durchlesen und ab hier!!!


%Figure~\ref{fig:trilemma} illustrates the circular structure of the pandemic trilemma.

%\begin{figure}[ht]
%\centering
%\begin{tikzpicture}[
    %node distance=3.5cm,
    %objective/.style={draw, rounded corners, minimum width=3.2cm, minimum height=1.2cm, align=center, font=\small\bfseries, line width=1pt},
    %instrument/.style={draw, dashed, rounded corners, minimum width=2.5cm, minimum height=0.9cm, align=center, font=\small, fill=gray!10},
    %arr/.style={-{Stealth[length=3mm]}, thick},
    %tension/.style={-{Stealth[length=3mm]}, thick, red!70!black, densely dashed},
    %label/.style={font=\scriptsize, fill=white, inner sep=2pt}
%]

% Objectives (triangle)
%\node[objective, fill=red!15] (health) at (90:4.2cm) {Public Health \\ $d_k \to 0$};
%\node[objective, fill=blue!15] (econ) at (210:4.2cm) {Economic Stability \\ $y_k \to 0$};
%\node[objective, fill=orange!15] (fiscal) at (330:4.2cm) {Fiscal Sustainability \\ $b_k \to 0$};

% Instruments
%\node[instrument] (S) at ($(health)!0.5!(econ) + (-1.8,0)$) {$\Sk$ (NPI)};
%\node[instrument] (F) at ($(econ)!0.5!(fiscal) + (0,-0.8)$) {$\Fk$ (Fiscal)};

% Positive effects (green)
%\draw[arr, green!50!black] (S) -- node[label, left, pos=0.4] {$-\phi_S$} (health);
%\draw[arr, green!50!black] (F) -- node[label, below left, pos=0.4] {$+\alpha_F$} (econ);

% Negative effects (red dashed)
%\draw[tension] (S) -- node[label, right, pos=0.4] {$-\alpha_S$} (econ);
%\draw[tension] (F) -- node[label, below right, pos=0.4] {$+\kappa_F$} (fiscal);

% Automatic stabiliser feedback
%\draw[tension, bend right=25] (econ) to node[label, below, pos=0.5] {$-\gamma_y$ (auto.\ stab.)} (fiscal);

% Explosive dynamic
%\draw[arr, red!70!black, bend left=30, densely dotted, thick] (health.east) to node[label, right, pos=0.5] {$\rho_\theta > 1$ if $\Sk = 0$} (health.north east);

%\end{tikzpicture}
%\caption{The circular structure of the pandemic trilemma. Solid green arrows denote positive policy effects; dashed red arrows denote adverse side effects. The suppression instrument $\Sk$ protects health but contracts the economy (Properties~1--2); the compensation instrument $\Fk$ stabilises output but accumulates debt (Property~3); automatic stabilisers transmit economic contraction to the fiscal position independently of active fiscal policy (Property~4). The explosive epidemiological dynamic ($\rho_\theta > 1$) ensures that inaction is not a viable strategy (Property~1).}
%\label{fig:trilemma}
%\end{figure}


%%Old Model and Literature Review



%%OOOLD
\section{Model Framework}

\subsection{Specifications}
\subsubsection{State and Control Vectors}

Following the notation of \citet{bertsekas2017dynamic}, the state of the system at time $k$  is represented by the vector $\bm{x}_k \in \mathbb{R}^4$:


\begin{itemize}
    \item $y_k$: \textbf{Output Gap} — The percentage deviation of real GDP from its potential trend (e.g., $-0.01$ represents a 1\% output loss).
    \item $d_k$: \textbf{Excess Mortality} — The deviation of the periodic death rate from the historical seasonal average, normalized by population.
    \item $b_k$: \textbf{Public Debt} — The deviation of the debt-to-GDP ratio from the pre-pandemic trend level, measured in percentage points.
    \item $\theta_k$: \textbf{Infection Rate} — The viral prevalence in the population, expressed as the share of actively infected individuals (where $\theta_k=0$ denotes eradication).
\end{itemize}

\noindent The planner's policy tools are represented by $\bm{u}_k \in \mathbb{R}^2$:
\begin{equation}
\bm{u}_k = \begin{pmatrix} S_k & F_k \end{pmatrix}'
\end{equation}
where $S_k$ denotes \textbf{Social Distancing} intensity (non-pharmaceutical interventions) and $F_k$ denotes \textbf{Fiscal Stimulus} intensity.


\begin{table}[ht]
\centering
\caption{Variables and Parameters Introduced in the Trilemma Analysis}
\label{tab:notation}
\renewcommand{\arraystretch}{1.25}
\small
\begin{tabular}{@{} c l l c @{}}
\toprule
\textbf{Symbol} & \textbf{Description} & \textbf{Type} & \textbf{Property} \\
\midrule
$\theta_k$ & Infection rate (viral prevalence) & State variable & 1 \\
$d_k$ & Excess mortality & State variable & 1 \\
$y_k$ & Output gap (\% deviation from potential) & State variable & 2 \\
$b_k$ & Public debt (deviation of debt-to-GDP ratio) & State variable & 3 \\
\midrule
$\Sk$ & Social distancing intensity (NPI) & Control variable & 2 \\
$\Fk$ & Fiscal stimulus intensity & Control variable & 3 \\
\midrule
$\rho_\theta$ & Epidemiological persistence ($>1$: explosive) & Parameter & 1 \\
$\delta_\theta$ & Infection-to-mortality transmission & Parameter & 1 \\
$\varepsilon_k$ & Exogenous epidemiological shock & Disturbance & 1 \\
$\alpha_S$ & Contractionary effect of NPI on output & Parameter & 2 \\
$\phi_S$ & Suppression effect of NPI on infections & Parameter & 2 \\
$\alpha_F$ & Fiscal multiplier (stimulus $\to$ output) & Parameter & 3 \\
$\kappa_F$ & Budgetary cost of fiscal stimulus & Parameter & 3 \\
$r$ & Real interest rate on sovereign debt & Parameter & 3 \\
$\gamma_y$ & Budget semi-elasticity (auto.\ stabilisers) & Parameter & 4 \\
\bottomrule
\end{tabular}
\end{table}





\vspace{2mm}

\noindent \textit{\textbf{Disclaimer}: As this is a first draft to start the model at a simple baseline, there could be a couple of extensions and refinements. I remarked those at the corresponding space in \textcolor{blue}{blue}.}

\subsubsection{System Dynamics and State-Space Representation}

To formalise the social planner’s trilemma, the system dynamics are defined as a discrete-time linear state-space model. The system evolves according to the following transition law:
\begin{equation}
    \bm{x}_{k+1} = \bm{A} \bm{x}_k + \bm{B} \bm{u}_k +  G\epsilon_k
\end{equation}

\noindent All state variables in the vector $\bm{{x}_k}' = (y_k, d_k, b_k, \theta_k)'$ are defined as deviations from the pre-pandemic steady state. Consequently, the origin of the state space represents the desired equilibrium to which the planner seeks to return the system. The transition is mapped as follows:

\begin{equation}
\begin{pmatrix} y_{k+1} \\ d_{k+1} \\ b_{k+1} \\ \theta_{k+1} \end{pmatrix} = 
\underset{\underbrace{\rule{1.2cm}{0pt}}_{\bm{A}}}{
\begin{pmatrix} 
\rho_y & 0 & 0 & 0 \\ 
0 & 0 & 0 & \delta_\theta \\ 
-\gamma_y & 0 & (1+r) & 0 \\ 
0 & 0 & 0 & \rho_\theta 
\end{pmatrix}}
\begin{pmatrix} y_k \\ d_k \\ b_k \\ \theta_k \end{pmatrix} +
\underset{\underbrace{\rule{0.8cm}{0pt}}_{\bm{B}}}{
\begin{pmatrix} 
-\alpha_s & \alpha_F \\ 
0 & 0 \\ 
0 & \kappa_F \\ 
-\phi_s & 0 
\end{pmatrix}}
\begin{pmatrix} S_k \\ F_k \end{pmatrix} + 
\underset{\underbrace{\rule{0.5cm}{0pt}}_{\bm{G}}}{
\begin{pmatrix} 0 \\ 0 \\ 0 \\ 1 \end{pmatrix}} 
\varepsilon_k
\end{equation}

\paragraph{Equation I: Output Gap Dynamics ($y_{k+1}$)}The first row of the system defines the evolution of the economic state:
\begin{equation}y_{k+1} = \rho_y y_k - \alpha_s S_k + \alpha_F F_k\end{equation}

The economic output gap is primarily driven by its own persistence ($\rho_y$), the contractionary effect of NPI stringency ($-\alpha_s$), and the stimulative effect of fiscal support ($\alpha_F$). We assume that the death rate ($d_k$), the fiscal deviation ($b_k$), and the viral prevalence ($\theta_k$) do not exert a \textit{direct} structural impact on the output gap of the subsequent period without an intervention of the social planner. \textit{\textcolor{blue}{While a high viral load or mortality may indirectly affect the economy through either behavioral changes "Fear Effect" or lower productive human capital stock, this model attributes such effects in a first step only to the social planner's response ($S_k$), maintaining the focus on policy-driven trade-offs first. To better capture the heterogeneity of fiscal interventions, a further objective is to refine $F_k$ by distinguishing between revenue- and expenditure-based policies and their respective targets.Additionally, an error term could be introduced to represent stochastic shocks from non-governmental activities. However, due to the Certainty Equivalence (CE) property, such shocks would not change the optimal policy rules derived in this model.}}


\paragraph{Equation II: Mortality Dynamics ($d_{k+1}$)} 
The second row of the system governs the evolution of excess mortality:
\begin{equation}
    d_{k+2} = \delta_\theta \cdot \underbrace{\bigl[\rho_\theta (1 - \phi_S S_k) \, \theta_k + \varepsilon_k\bigr]}_{\theta_{k+1}} 
    \label{eq:excess dynamics}
\end{equation}
In this specification, mortality is solely driven by the lagged viral prevalence ($\delta_\theta$). This choice is rooted in \textit{biological realism}: pandemic controls are primarily designed to suppress viral transmission ($\theta$) rather than directly altering the clinical outcome of individuals already infected. By channeling the effect of $S_k$ through the epidemiological state ($\theta$), we introduce a structural time lag between policy intervention and its impact on mortality. Empirical evidence suggests that the delay between infection (cases) and death averages a lag of eight weeks \citep{Gagnon2023}, meaning that mortality serves as a trailing indicator of the pandemic's trajectory. \textit{\textcolor{blue}{The only structural argument for incorporating $S_k$ directly into the mortality equation would be interventions specifically targeting the protection of vulnerable groups, such as the localized shielding of \textit{nursing homes}, which could decouple mortality from general prevalence.}}

 
 
\paragraph{Equation III: Fiscal Dynamics ($b_{k+1}$)} 
The third row represents the intertemporal budget constraint of the social planner, governing the evolution of public debt:
\begin{equation}
    b_{k+1} = (1+r) b_k - \gamma_y y_k + \kappa_F F_k
\end{equation}
The debt level, defined as the deviation from the pre-pandemic trend, evolves based on the real interest rate ($r$), reflecting the burden of debt servicing. The debt level is also augmented by automatic stabilizers that captures the endogenous link between economic performance and the budget. A negative output gap leads to a decline in tax revenues—such as Value Added Tax (VAT) or income tax—while simultaneously increasing ordinary mandatory expenditures. Conversely, economic recovery (a narrower output gap) inherently alleviates fiscal pressure through for example higher tax receipts . Consequently, $\gamma_y$ acts as a sensitivity parameter that aggregates all non-discretionary budgetary responses to economic fluctuations. This term also implicitly captures the fiscal legacy of past interventions; for instance, if previous stimulus measures successfully narrowed the output gap, the resulting fiscal pressure through automatic stabilizers is reduced. Third, the term $\kappa_F F_k$ represents the direct budgetary impact of the planner’s fiscal support programs. Where it is in this first version only represented by the parameter $\kappa_F$ and their aggregate impact on sovereign debt. Notably, this specification assumes no direct enforcement costs of NPIs ($S_k$) on the budget ($b_{k+1}$), as these are considered to be part of $F_k$ (e.g., health spending) or negligible compared to the massive indirect impact of lockdowns via the output gap. This emphasizes that the fiscal cost of social distancing is primarily an opportunity cost of lost economic activity rather than a direct administrative expense. \textit{\textcolor{blue}{Possible extensions may disaggregate $F$ by design (expenditure vs. revenue) and target (households vs. firms) to capture policy heterogeneity. Additionally, incorporating a crowding-out effect ($b \to y$) would allow the model to account for the intertemporal feedback of debt on the output gap.}}


\paragraph{Equation IV: Epidemiological Dynamics ($\theta_{k+1}$)}The final row describes the progression of the virus:\begin{equation}\theta_{k+1} = \rho_\theta \theta_k - \phi_s S_k + \varepsilon_k\end{equation}The pandemic follows a purely biological persistence ($\rho_\theta$), modulated only by the intensity of NPIs ($-\phi_s$) and an exogenous stochastic shocks ($\varepsilon_k$). We assume that viral spread is independent of the output gap ($y_k$) or the level of public debt ($b_k$). Behavioural adjustments by the populace are assumed to be a function of mandated NPIs ($S_k$) as well as pandemic pressure. Note that the health state $d_k$ is not directly controllable; the planner can influence mortality only indirectly through suppressing 
viral prevalence, with an inherent policy lag of at least one period. \textit{ \textcolor{blue}{All these elements can be recognized as in the error term for the time being.}}

\vspace{4mm}

\noindent Further assumptions that follows the framework of \cite{bertsekas2017dynamic} (Chapter 3.1), is that we assume that $\varepsilon_k$ is a white noise process with $E[\varepsilon_k] = 0$ and $E[\varepsilon_k^2] = \sigma^2$, representing exogenous shocks such as viral mutations. Because the system is linear and the cost function is quadratic, the Certainty Equivalence Principle applies. This ensures that the optimal feedback gain $\bm{K}_k$ is identical to the deterministic case, although the expected minimum cost will include an additional term proportional to the noise variance $\sigma^2$.


%MACHEN

\subsubsection{Control Leverage and Controllability}


The structure of the control matrix $\bm{B}$ reveals a fundamental asymmetry in the planner's policy instruments. Social distancing ($S_k$) operates as a instrument that simultaneously affecting output gap and viral prevalence. Fiscal stimulus ($F_k$) functions as the second instrument targeting the output gap while directly augmenting public debt. This complementarity is essential; since $rank(\bm{B}) < dim(\bm{x})$, neither instrument alone can span the full state space, necessitating a coordinated policy mix. 

A critical feature of this system is that mortality ($d_k$) cannot be directly controlled, as the second row of $\bm{B}$ consists of zeros. The planner influences health outcomes only indirectly through the transmission channel $S_k \to \theta_{k+1} \to d_{k+2}$. This structure embeds an irreducible policy lag: interventions enacted in period $k$ affect mortality no earlier than period $k+2$. 

\vspace{2mm}

\noindent Despite this constraint, the system remains fully controllable. The controllability matrix

\begin{equation}
    \mathcal{C} = \begin{pmatrix} \bm{B} & \bm{AB} & \bm{A}^2\bm{B} & \bm{A}^3\bm{B} \end{pmatrix}
\end{equation}
maintains full row rank under non-degeneracy conditions ($\delta_\theta, \phi_s, \alpha_s, \alpha_F, \kappa_F \neq 0$). While $d_k$ is decoupled from $\bm{B}$, it is reachable through $\bm{AB}$ via the interaction term $\delta_\theta \cdot \phi_s$. 

This controllability structure carries two major implications. First, the health-economy trade-off is temporally asymmetric: fiscal stimulus stabilizes output contemporaneously, whereas the mortality benefits of social distancing materialize with a multi-period delay. Second, the planner faces a \textit{minimum-time} constraint on health recovery that no policy intensity can circumvent—a feature that mathematically rationalizes early, preemptive intervention when pandemic trajectories are still manageable.

The eigenvalues of $\bm{A}$ are $\{\rho_y, 0, (1+r), \rho_\theta\}$. Under standard parameterizations ($\rho_y, \rho_\theta \in (0,1)$ and $r > 0$), the debt state is explosive in the absence of intervention: $(1+r) > 1$ lies outside the unit circle. This instability underscores 
that the planner's problem is not merely optimization but stabilization—active policy is required to return the system to its pre-pandemic equilibrium.

\vspace{5mm}

\textit{Remark: While $F_k$ is currently treated as an aggregate, a further disaggregation of fiscal policy by design (expenditure vs. revenue) and target (households vs. firms) offers a compelling path for future refinement. Such an extension would significantly increase the strategic depth of the model, as it allows for the analysis of heterogeneous fiscal multipliers and their specific trade-offs between immediate economic stabilization and long-term debt sustainability, making the planner's optimization problem even more nuanced.}

\vspace{2mm}



\section{The Optimization Problem}

\subsection{Objective Function}
The social planner seeks to minimize the expected intertemporal cost:
\begin{equation}
J = \mathbb{E} \left\{ \bm{x}_N' \bm{Q}_N \bm{x}_N + \sum_{k=0}^{N-1} \left( \bm{x}_k' \bm{Q}_k \bm{x}_k + \bm{u}_k' \bm{R} \bm{u}_k \right) \right\}
\end{equation}
where $\bm{Q}_k$ penalizes deviations of state variables from the 
pre-pandemic equilibrium, $\bm{Q}_N$ imposes a terminal cost, and 
$\bm{R}$ reflects the direct cost of policy implementation.

\subsection{Weighting Matrices}

\subsubsection{Running State Cost (economical costs) ($\bm{Q}_k$)}
\begin{equation}
\bm{Q}_k = \beta^k 
\begin{pmatrix} 
w_y & 0 & 0 & 0 \\ 
0 & w_d & 0 & 0 \\ 
0 & 0 & w_b & 0 \\ 
0 & 0 & 0 & 0 
\end{pmatrix}
\end{equation}
The discount factor $\beta \in (0,1)$ reflects the planner's time 
preference. The zero weight on $\theta_k$ implies that viral 
prevalence is not penalized directly; rather, infections enter the 
objective only through their delayed impact on mortality via the 
transition $d_{k+1} = \delta_\theta \theta_k$.

\subsubsection{Terminal State Cost ($\bm{Q}_N$)}
\begin{equation}
\bm{Q}_N = \beta^N
\begin{pmatrix} 
0 & 0 & 0 & 0 \\ 
0 & 0 & 0 & 0 \\ 
0 & 0 & W_b & 0 \\ 
0 & 0 & 0 & 0 
\end{pmatrix}
\end{equation}
The terminal cost focuses exclusively on public debt. Output gaps 
close endogenously as restrictions lift; excess deaths are sunk 
costs already penalized through $w_d$. Debt, by contrast, compounds 
at rate $(1+r)$ and persists beyond the planning horizon. The 
formulation of $W_b$ captures the present value of this 
accumulation.

The post-COVID period underscores this concern: inflation surges 
prompted sharp policy rate increases, raising sovereign borrowing 
costs across advanced economies \citep{JORDA2023104474}. Higher $r$ 
tightens the convergence condition and amplifies $W_b$, reinforcing 
that fiscal prudence during the crisis is a precondition for 
post-pandemic stability.

\subsubsection{Control Cost (political costs) ($\bm{R}$)}
\begin{equation}
\bm{R} = 
\begin{pmatrix} 
r_S & 0 \\ 
0 & r_F 
\end{pmatrix}
\end{equation}
The parameters $r_S$ and $r_F$ represent the direct costs of 
deploying social distancing measures and fiscal stimulus, 
respectively. These costs are independent of the state variables and 
capture political, administrative, or welfare costs of intervention 
per se.

\subsection{Mathematical Requirements}
For the LQR problem to admit a unique solution, the following 
conditions must hold:
\begin{itemize}
    \item \textbf{Positive semidefiniteness} ($\bm{Q}_k, \bm{Q}_N 
    \succeq 0$): Satisfied if $w_y, w_d, w_b, W_b \geq 0$. 
    This ensures that state deviations are penalized, not rewarded.
    \item \textbf{Positive definiteness} ($\bm{R} \succ 0$): Satisfied 
    if $r_S, r_F > 0$. This guarantees a unique optimum by ensuring 
    that policy interventions are costly at the margin, precluding 
    unbounded control.
\end{itemize}


\subsection{Solution via Dynamic Programming}
The optimal control problem admits a closed-form solution via 
backward induction (see Appendix) \citep[Chapter 3.1]{bertsekas2017dynamic}. 
The optimal policy is a linear state-feedback law:
\begin{equation}
    \bm{u}_k^* = -\bm{K}_k \bm{x}_k
\end{equation}
where the gain matrix $\bm{K}_k$ and value function matrix 
$\bm{P}_k$ are computed recursively from the terminal condition 
$\bm{P}_N = \bm{Q}_N$:
\begin{align}
    \bm{K}_k &= \left( \bm{B}' \bm{P}_{k+1} \bm{B} + \bm{R} \right)^{-1} 
    \bm{B}' \bm{P}_{k+1} \bm{A} \\[1ex]
    \bm{P}_k &= \bm{A}' \bm{P}_{k+1} \bm{A} - \bm{A}' \bm{P}_{k+1} \bm{B} 
    \, \bm{K}_k + \bm{Q}_k
\end{align}
The recursion is solved numerically for $k = N-1, \ldots, 0$.


\subsection{Certainty Equivalence}
The stochastic disturbance $w_k$ enters linearly and the cost function 
is quadratic, satisfying the conditions for certainty equivalence. 
Consequently:
\begin{enumerate}
    \item The optimal feedback gain $\bm{K}_k$ is identical to the 
    deterministic case ($w_k \equiv 0$).
    \item The expected cost increases by 
    $\sum_{k=0}^{N-1} \mathbb{E}[w_k' \bm{P}_{k+1} w_k]$, reflecting 
    the irreducible cost of pandemic uncertainty.
    \item The planner optimizes as if future shocks equal their 
    expected value (zero), then absorbs realized deviations through 
    the feedback structure.
\end{enumerate}

\section{Literature Review}

%Für Erkenntnisse während COVID erst später bei der kalbirierung vom modell-> Literatur Review
The pandemic environment introduces structural features that modify the transmission mechanisms of all three instrument categories relative to conventional recessions.

The most fundamental complication is that fiscal stimulus must operate within an economy whose absorptive capacity is simultaneously constrained by the suppression measures that necessitate the stimulus. When NPIs restrict the supply of goods and services available for purchase, demand-side fiscal impulses cannot translate into consumption at the rate observed under normal conditions. This mechanism has been documented empirically: \cite{coibion2020} find that only approximately 40 percent of US stimulus payments were spent, with the remainder saved, and that spending propensities were substantially higher among unemployed and financially constrained recipients. \cite{chetty2020} confirm that stimulus checks increased expenditure among low-income households but not in sectors directly affected by containment restrictions. \cite{baker2020} report that approximately one-third of stimulus payments were spent within the first weeks, with significantly larger effects among lower-income recipients.

These findings suggest that $\alpha_F^j$ under pandemic conditions is not a structural constant but is \textit{state-dependent}---conditioned on the prevailing level of suppression $S_k$. \cite{auerbach2022} provide the most direct evidence for this state-dependence: exploiting variation in US Department of Defense spending across cities with and without lockdown restrictions, they find that government spending increased employment only in cities without meaningful containment measures. In restricted cities, the employment effect was statistically indistinguishable from zero. This result implies that the fiscal multiplier can collapse toward zero under strict containment, fundamentally altering the planner's trade-off calculus.

Beyond the suppression of demand-side transmission, \cite{ghassibe2021} demonstrate theoretically that demand-stimulating measures are less effective in supply-induced downturns more generally---a finding that extends beyond the pandemic context but applies with particular force to an environment where the economic contraction originates from an exogenous restriction on production and exchange rather than from an endogenous demand shortfall.

The pandemic context also reveals the potential importance of credit instruments. While loans and guarantees exhibit limited \textit{direct} demand effects, their role in preserving productive capacity may generate substantial \textit{indirect} output benefits during a temporary crisis. \cite{gourinchas2021} find that fiscal support approximately halved the increase in SME insolvencies, preventing permanent destruction of firm-specific capital and employment relationships. Similarly, European job retention schemes preserved labor market attachment more effectively than the US approach of unemployment insurance combined with rehiring, resulting in faster labor market recovery upon reopening \citep{aum2021}. These findings suggest that the conventional multiplier framework may understate the output contribution of credit and preservation-oriented instruments during a temporary, supply-constrained downturn.

Cross-country evidence further supports the aggregate effectiveness of fiscal support during the pandemic, while leaving open the question of instrument-specific contributions. \cite{chudik2021} employ a threshold-augmented Global VAR model covering 33 countries and find that countries with larger fiscal packages experienced significantly smaller output contractions. The \cite{imf2022fiscal} estimates that cumulative fiscal multipliers one year after a health crisis are approximately twice as large as during normal recessions, reflecting pent-up demand effects as containment is relaxed. However, both of these estimates operate at the aggregate level of total fiscal support and do not distinguish between the contributions of individual instrument categories.

For expenditure-based measures, the pre-pandemic consensus places government spending multipliers in the range of approximately 0.5 to 1.5, with substantial variation across estimation methods and economic contexts. \cite{ramey2018} survey the literature and find that multipliers are generally below unity in normal times but rise at the zero lower bound of monetary policy, where monetary accommodation amplifies fiscal transmission. Cross-sectional estimates exploiting geographic variation in US government spending during the American Recovery and Reinvestment Act suggest multipliers at the upper end of this range \citep{ARRA_QUELLE}. Within expenditure-based instruments, the multiplier depends on the recipient's marginal propensity to consume (MPC): unemployment insurance benefits, which reach liquidity-constrained households with high MPCs, tend to exhibit larger output effects than untargeted lump-sum transfers \citep{fariae2020, bayer2020}.

For revenue-based measures, the evidence consistently indicates smaller output effects than expenditure-based instruments. The transmission from tax relief to aggregate demand is attenuated by the saving behavior of recipients, which interposes a behavioral wedge between the fiscal impulse and the demand response. The extent of this attenuation depends on the type of tax instrument: temporary reductions in consumption taxes (such as VAT cuts) may induce intertemporal substitution of expenditure, while income tax reductions operate through the slower channel of disposable income.

For loans and guarantees, direct output multiplier estimates are sparse in the pre-pandemic literature. The primary mechanism---liquidity provision and insolvency prevention---does not map neatly onto the demand-side multiplier framework that applies to expenditure and revenue measures. Credit support affects output primarily by preventing \textit{negative} outcomes (firm closure, job destruction) rather than by generating \textit{positive} demand impulses, making its contribution to $\alpha_F^{L\&G}$ qualitatively distinct and difficult to estimate using standard multiplier methodologies.



Within this environment, the existing literature has addressed the bilateral
components of the trilemma in isolation.

\paragraph{Containment and output.} The contractionary effect of NPIs on
economic activity is well established. Cross-country evidence documents GDP
reductions of 20--25\% in several OECD economies during initial lockdowns,
with emerging markets experiencing larger contractions for comparable
stringency levels due to higher employment shares in contact-intensive
sectors and weaker automatic stabilizers \citep{imf2020, portes2020}. The
temporal pattern is characteristic: restrictions contract output
contemporaneously, followed by a partial catch-up upon relaxation---consistent
with the autoregressive structure captured by $\rho_y$ in the theoretical
model. The timing of intervention matters: early, stringent, but time-limited
containment dominates delayed intervention \citep{acemoglu2021, alvarez2021,
balmford2020, kargi2023}, though optimal stringency depends heavily on the
assumed value of a statistical life \citep{alvarez2021}.

\paragraph{Containment and mortality.} The protective effect of lockdowns
operates through an epidemiological transmission channel with inherent policy
lags. Synthetic control estimates for Sweden suggest that approximately 9.5\%
of total mortality during February--September 2020 could have been avoided
under alternative containment policies \citep{born2021}. The stringency index
emerges as the most significant predictor of both infection rates and
mortality outcomes across specifications \citep{kargi2023}. Critically, the
lag between intervention and health outcomes---estimated at approximately
eight weeks from infection to death \citep{gagnon2023}---implies that
mortality serves as a trailing indicator, consistent with the two-period
transmission $S_k \to \theta_{k+1} \to d_{k+2}$ embedded in the model.

\paragraph{Fiscal support and output.} Pre-pandemic evidence placed
government expenditure multipliers at approximately 1.5
\citep{ramey2018}, but the pandemic context introduced
novel complications. Cumulative multipliers during health crises are
approximately twice as large as during normal recessions \citep{deb2021},
reflecting substantial economic slack. However, this aggregate finding masks
critical instrument heterogeneity. Broad-based household transfers yielded
compressed multipliers as recipients saved the majority of stimulus payments
under pandemic consumption constraints \citep{chetty2020, coibion2020}; the
U.S.\ Paycheck Protection Program illustrates the consequences of poor
targeting, with approximately three-quarters of \$800 billion in
disbursements accruing to top-quintile households \citep{autor2022}. The
transfer multiplier estimated by \citet{bayer2023} confirms that the
composition of fiscal packages---not merely their size---determines
stabilization outcomes. Firm-targeted support, particularly job retention
schemes and payroll subsidies, preserved employment relationships and
facilitated rapid recovery upon reopening \citep{bennedsen2020}.

\paragraph{State-dependency of fiscal effectiveness.} The interaction between
containment stringency and fiscal transmission is a central but
underexplored channel. \citet{auerbach2022} exploit variation in U.S.\
Department of Defense spending across cities with and without stay-at-home
orders during early 2020: government spending increased employment only in
unrestricted cities at a cost of approximately \$4,000 per job per month,
while the employment effect was statistically indistinguishable from zero
under binding containment. Notably, even in unrestricted cities, spending did
not stimulate local consumption---the standard high-MPC transmission channel
was inoperative, and employment effects operated through direct labor
absorption rather than induced demand.

\paragraph{Capacity preservation and firm survival.} Despite the deepest
recession in decades, OECD bankruptcy filings declined by 17\% in 2020
relative to 2019 \citep{djankov2021}, with an estimated 25,000 missing
insolvencies in Germany alone attributable to government support programs
\citep{banerjee2021}. \citet{gourinchas2021} provide the theoretical
rationale: credit support preserved viable firms in ``hibernation,''
maintaining stakeholder relationships---employer-employee matches, supplier
networks, customer bases---that would have been irreversibly destroyed in
insolvency and costly to rebuild upon reopening.

\paragraph{Fiscal support and debt.} The fiscal response generated
unprecedented debt increases: EU government debt rose by 8.4 percentage
points to 88\% of GDP within a single quarter, reaching 94\% by year-end
2020 \citep{eurostat2021}. The debt impact varies by instrument type:
above-the-line measures generate immediate deficit increases, while
guarantees create contingent liabilities with delayed and partial
realization. Critically, high existing debt levels blunt the impact of
additional stimulus through crowding out and expectations of future
consolidation, with multiplier estimates of 0.6--1.0 in low-debt countries
but below 0.4 in high-debt countries \citep{ilzetzki2013}.

\paragraph{The gap.} What remains unaddressed is the \textit{joint}
optimization across these bilateral channels: which fiscal composition---across
demand injection and capacity preservation, at what intensity, and at which
point in the containment trajectory---minimizes the cumulative debt footprint
for a given level of output stabilization and health protection?





\paragraph{The research gap.}

The preceding discussion reveals a tension in the existing literature. On one hand, the pre-pandemic evidence on fiscal multipliers provides well-established priors for the relative effectiveness of expenditure, revenue, and credit instruments under normal macroeconomic conditions. On the other hand, the pandemic introduces structural features---supply-constrained consumption, state-dependent transmission, the preservation function of credit instruments---that may modify these priors in ways that have not been systematically quantified.

The gap is threefold. First, the \textit{cross-instrument comparison} under pandemic conditions remains incomplete. Studies have examined individual instruments in isolation---stimulus checks \citep{coibion2020, chetty2020, baker2020}, payroll subsidies \citep{aum2021}, loan programs \citep{autor2022, gourinchas2021}---but no cross-country analysis has simultaneously estimated the output multipliers $\alpha_F^{Exp}$, $\alpha_F^{Rev}$, and $\alpha_F^{L\&G}$ within a unified framework that permits direct comparison.

Second, the \textit{interaction between fiscal effectiveness and containment stringency}---the state-dependence of $\alpha_F^j$ with respect to $S_k$---has been documented for aggregate government spending \citep{auerbach2022} but not for disaggregated instrument categories. Whether expenditure, revenue, and credit measures are equally suppressed under strict containment, or whether some instruments are more robust to supply-side constraints than others, is an open question with direct implications for optimal fiscal composition during a pandemic.

Third, the existing literature has largely treated fiscal multipliers and pandemic control as \textit{separate} domains. Studies on optimal lockdown policy typically abstract from the fiscal dimension \citep{alvarez2021, acemoglu2021}, while studies on pandemic fiscal responses rarely incorporate the endogenous interaction with containment stringency. The joint estimation of containment and fiscal effects---including their interaction---constitutes a central objective of the empirical analysis in Section~X.














